problem:  
Let \( \mathfrak{n} \) be a nilpotent Lie algebra equipped with an inner product \( \langle \cdot,\cdot\rangle \).  
Denote by \( \mathcal{V}(\mathfrak{n}) \) the algebraic variety of nilpotent Lie brackets on the fixed underlying vector space of \( \mathfrak{n} \).  A bracket \( [\cdot,\cdot]_0 \in \mathcal{V}(\mathfrak{n}) \) is called a **nilsoliton bracket** if there exists a constant \( c\in\mathbb R \) and a symmetric derivation \(D_E\) such that
\[
\mathrm{Ric}_{[\cdot,\cdot]_0} = c\,\mathrm{Id} + D_E .
\]
Heber and Lauret showed that such a pair \(([\cdot,\cdot]_0, D_E)\) algebraically determines a unique solvable Einstein Lie algebra \( \mathfrak{s}_0 = \mathbb R A \oplus \mathfrak{n} \) with
\[
[A,X] = D_E X, \qquad X\in\mathfrak{n}.
\]
Let \( \mathcal{E}_{\mathrm{alg}} \) denote the **algebraic extension operator** sending a nilsoliton bracket \( [\cdot,\cdot] \) with corresponding symmetric derivation \(D([\cdot,\cdot])\) to the solvable Einstein bracket
\[
\mathcal{E}_{\mathrm{alg}}([\cdot,\cdot])\;:=\;\big([\cdot,\cdot]_E,\ D([\cdot,\cdot])\big),
\]
where \( [A,X]_E = D([\cdot,\cdot])X \) and \( [X,Y]_E = [X,Y] \).

Now fix a nilsoliton bracket \( [\cdot,\cdot]_0 \) producing the Einstein bracket \( \mathcal{E}_{\mathrm{alg}}([\cdot,\cdot]_0) \).  Regard \( \mathcal{V}(\mathfrak{n}) \) as a subset of a finite-dimensional vector space (via structure constants).  The tangent space \( T_{[\cdot,\cdot]_0}\mathcal{V}(\mathfrak{n}) \) is the space of skew‑symmetric bilinear maps \( \varphi:\Lambda^2\mathfrak{n}\to\mathfrak{n} \) satisfying the Jacobi‑compatibility constraints.

> **Problem (Algebraic Local Injectivity of the Einstein Extension).**  
> Consider the linearization \( D\mathcal{E}_{\mathrm{alg}}|_{[\cdot,\cdot]_0} : T_{[\cdot,\cdot]_0}\mathcal{V}(\mathfrak{n}) \to T_{\mathcal{E}_{\mathrm{alg}}([\cdot,\cdot]_0)}\mathcal{V}(\mathfrak{s}_0) \).  
>  
> Is its kernel modulo the natural action of \( \mathrm{Der}(\mathfrak{n}) \) and scaling trivial?  
> Equivalently: does every infinitesimal deformation of a nilsoliton bracket that induces an isospectral infinitesimal deformation of its Einstein extension necessarily come from an infinitesimal automorphism or scaling of \( \mathfrak{n} \)?  

A positive answer would mean that the algebraic mechanism passing from nilsolitons to Einstein solvmanifolds is **locally injective** at the level of Lie brackets—distinct nearby nilsoliton deformations cannot yield different Einstein extensions infinitesimally indistinguishable. A negative answer would reveal new algebraic degeneracies between Ricci soliton and Einstein moduli.

Why is it a "good" problem:  
1. **Mathematical precision:**  
   The reformulation replaces ill‑defined moduli spaces with the explicit **variety of Lie brackets**, a classical algebraic object admitting a well‑defined tangent‑space structure. The linearization \(D\mathcal{E}_{\mathrm{alg}}\) is now concretely interpretable as a derivation‑dependent linear map between finite‑dimensional spaces of structure tensors—making the question rigorous.

2. **Profound geometric insight:**  
   Determining whether \(D\mathcal{E}_{\mathrm{alg}}\) has trivial kernel (mod automorphisms) tests an infinitesimal correspondence between two geometric fixed‑point conditions:  
   *the Ricci soliton equation* on \(\mathfrak{n}\) and *the Einstein condition* on its solvable extension.  
   It thus probes how rigidity in the Ricci soliton realm transfers—or fails to transfer—to Einstein solvmanifolds.

3. **Bridge between fields:**  
   It connects **Lie algebraic deformation theory** (varieties of brackets), **Ricci soliton analysis**, and **Einstein metric geometry**. Solving it may require combining algebraic geometry (local behavior of group orbits), analysis (Ricci operator differentials), and Riemannian geometry (Einstein deformation spaces).

4. **Extreme simplicity, deep impact:**  
   The question—whether any infinitesimal nilsoliton deformation disappears under Einstein extension—is compact and transparent. Yet affirmative resolution would underwrite a new algebraic rigidity principle; a counterexample would expose a subtle failure in the Lauret–Heber correspondence and enrich the taxonomy of homogeneous Einstein spaces.

Hence the problem is mathematically sound, conceptually central, and connects disparate geometric frameworks—meeting all defining traits of a “good” research‑level problem.